\section*{Appendix A: Revision History}

This framework was developed across six rounds of external review, each targeting a different weakness in the prior draft.

\begin{center}
\begin{tabular}{|p{0.08\textwidth}|p{0.20\textwidth}|>{\small\raggedright\arraybackslash}p{0.60\textwidth}|}
\hline
Round & Focus & What changed \\
\hline
1 & Evidentiary accuracy & Corrected overgeneralized claims, misattributed sources (Kahneman \& Klein, Cynefin, the Bitter Lesson, Parasuraman/Sheridan/Wickens), and a factual error in the Terminus-4B case study. Named and dated the OpenAI/Hugging Face containment incident precisely rather than leaving it anonymous. Softened claims about privilege waiver, output reproducibility, and model-confidence calibration that had been stated more strongly than the evidence supports. \\
\hline
2 & Systems safety & Identified a genuine structural gap: the framework had substantial coverage of where errors originate and why verification is hard, but no explicit account of how an erroneous or uncertain state propagates through a socio-technical system and becomes --- or fails to become --- real harm. Added the material now in Sections 11 through 14: hazard-first analysis, propagation topology, verification as an assurance argument, and operational resilience. \\
\hline
3 & Purpose and standard of success & Observed that a long catalogue of failure modes, without a stated positive objective, reads as an argument against ever deploying anything. Added Section 2, establishing comparative reliability against a strong human baseline --- not perfection --- as the framework's actual standard of success. \\
\hline
4 & Job definition & Found that job definition, on which every later section depends, had been compressed into a single sentence. Expanded Section 3 into a full treatment distinguishing the job, the decision it informs, the task assigned, the output produced, and the effect that follows --- and introduced job drift as a named failure mode. \\
\hline
5 & Navigability & Found the report dense enough that a senior, time-pressed reader had no way to tell which sections carry the actual decision versus which build and pressure-test it. Added a front-matter guide mapping the report's structure to the specific costly misdiagnosis each part exists to prevent, without shortening or simplifying the underlying material. \\
\hline
6 & De-formalization & Found that earlier revisions had drifted toward prescribing a governance process rather than supporting the judgment a decision-maker exercises. Replaced the four-outcome decision formalism in Section 3 with an open-ended standard for a defensible conclusion; renamed the use-and-reliance "contract" to a "description"; added an explicit proportionality principle so the report's own machinery scales with actual stakes rather than applying uniformly; moved the Revision History and the full set of open questions out of the main reading path into appendices; labeled Section 8 as dated context rather than enduring framework; and added a Closing Perspective distinguishing what the framework can establish from the judgment it still leaves to the reader. \\
\hline
7 & Stakeholder lens accuracy & A role-oriented (CTO/CISO/CRO/Executive Committee) summary was proposed as a companion document. Graded against the main text, it showed a consistent pattern: careful, bounded claims compressed into stronger, more absolute-sounding ones, including two near-verbatim regressions to overclaims Round 1 had specifically corrected (confidentiality leak vectors framed as routine rather than conditional; reviewer complacency framed with a specific, unmeasured trajectory). Added Appendix D with the corrected lens summary, each section closing with the open question most relevant to that role rather than presenting the material as fully settled. \\
\hline
\end{tabular}
\end{center}
